\chapter{Scheduler, memory mapping, interface}

\section{연산량보다 빈 cycle을 먼저 센다}

같은 butterfly 수라도 layer 시작, BRAM read, pipeline fill/drain, 다음 layer
descriptor 준비 때문에 cycle 수가 달라진다. 최종 radix-2 schedule의 관측식은
다음처럼 정리된다.
\[
C_{core}=L\,(N_{issue/layer}+15).
\]
ML-KEM은 7개 layer와 layer당 64 issue로 $7(64+15)=553$ cycle,
ML-DSA/HAETAE는 8개 layer로 $8(64+15)=632$ cycle이다. 이 식의 15는 단일
butterfly latency만 뜻하지 않는다. address issue부터 final writeback과 layer
전환까지의 관측 overhead다.

최종 FSM은 다음 layer descriptor를 미리 계산하고, 마지막 writeback edge의
lookahead로 다음 layer를 시작한다. address generator도 first transaction seed를
descriptor에 포함해 layer 시작 뒤의 별도 준비 cycle을 줄인다.

\section{NTRU+864의 half-empty issue 제거}

NTRU+864의 특수 step-3 구간을 일반 schedule로 표현하면 일부 transaction에서
lane을 충분히 사용하지 못한다. 최종 \rtlfile{ntt\_addrgen}은
\rtlfile{ADDR\_MODE\_N864\_STEP3}를 별도로 두어 유효 coefficient 조합을 직접
발행한다. 이는 memory layout 자체를 바꾸는 mode가 아니다. logical address
schedule만 바꾸고 affine mapping은 유지한다.

이 구분 덕분에 scheduler test는 issue address/valid를 검사하고, mapper test는
bijection/bank conflict를 독립적으로 검사할 수 있다. schedule과 physical layout을
한 always block에 섞으면 어느 쪽의 오류인지 분리하기 어렵다.

\section{Core-cycle 결과}

\begin{longtable}{L{31mm}r r r r}
\toprule
Configuration & 최초 cycle & 최종 cycle & Cycle 감소 & 100/150 MHz 시간 감소 \\
\midrule
ML-KEM-256 & 557 & 553 & 0.7\% & 33.8\% \\
ML-DSA-256 & 636 & 632 & 0.6\% & 33.8\% \\
HAETAE-256 & 636 & 632 & 0.6\% & 33.8\% \\
NTRU+768 & 1,709 & 1,513 & 11.5\% & 41.0\% \\
NTRU+864 & 2,053 & 1,761 & 14.2\% & 42.8\% \\
NTRU+1152 & 2,605 & 2,313 & 11.2\% & 40.8\% \\
\bottomrule
\end{longtable}

마지막 열은 최초 100~MHz와 최종 150~MHz implementation point를 사용한 core-only
시간이다. radix-2-only mode는 cycle 감소가 작고 주파수 상승 효과가 중심이다.
Mixed-radix NTRU+는 scheduler cycle 감소가 추가로 기여한다.

\section{AXI LOAD: beat 안의 병렬성을 memory lane으로 전달한다}

64-bit AXI input 한 beat는 ML-DSA에서 32-bit coefficient 두 개, narrow mode에서
16-bit coefficient 네 개를 담는다. 최종 loader는 accepted beat마다 최대 네
logical write를 같은 cycle에 만든다.

\begin{lstlisting}
oLOAD_WE  <= 1'b1;  oLOAD_ADDR  <= coeff_q;
oLOAD1_WE <= 1'b1;  oLOAD1_ADDR <= coeff_q + 1;
oLOAD2_WE <= !mldsa;oLOAD2_ADDR <= coeff_q + 2;
oLOAD3_WE <= !mldsa;oLOAD3_ADDR <= coeff_q + 3;
\end{lstlisting}

지원하는 모든 $N$이 packing factor로 정확히 나누어지므로 각 lane의 upper-bound
comparator를 반복할 필요가 없다. 한 번의 \rtlfile{all\_issued/final\_accept}
판정으로 충분하다. 불변조건을 이용해 세 개의 11-bit 비교기를 제거한 사례다.

\rtlfile{oWS\_TREADY}는 receive state에서 항상 1이고 내부 bubble을 삽입하지
않는다. 따라서 source가 TVALID를 연속으로 주면 매 cycle 한 beat를 받는다.

\section{AXI DUMP: backpressure보다 먼저 reservation한다}

BRAM request는 발행 후 취소할 수 없다. downstream \rtlfile{TREADY}가 내려간
뒤 FIFO가 꽉 차면 이미 날아오는 response를 버리게 된다. 최종 dump는 현재
FIFO count만 보지 않고, 아직 도착하지 않은 request까지 slot을 예약한다.

\begin{lstlisting}
wire pop       = (count_q != 0) && iWM_TREADY;
wire can_issue = (credit_q != 0) || pop;
wire issue     = active_q && !all_issued && can_issue;

case ({issue,pop})
  2'b10: credit_q <= credit_q - 1'b1;
  2'b01: credit_q <= credit_q + 1'b1;
endcase
\end{lstlisting}

같은 cycle에 pop이 있으면 그 slot을 새 request에 즉시 재사용한다. 정상
\rtlfile{TREADY=1} 상태에서 refill bubble 없이 \ii=1을 유지하면서, 임의의
backpressure에서는 overflow를 막는다. FIFO payload는 reset하지 않는다.
\rtlfile{count\_q=0}이 모든 오래된 entry를 invalid로 만들기 때문에 payload reset을
없애 LUTRAM inference와 packing을 돕는다.

\section{Interface 두 개를 비교하는 법}

AXI4-Stream과 AXI BRAM Controller는 같은 engine을 사용하지만 PS--PL data movement가
다르다. 따라서 비교에는 두 층이 있다.

\begin{itemize}
  \item \textbf{Core-only}: 같은 engine의 cycle와 standalone area
  \item \textbf{End-to-end}: configuration, cache handling, DMA/MMIO transfer,
        polling, output availability를 포함한 board 시간
\end{itemize}

최종 complete-system에서 stream 쪽은 AXI DMA MM2S/S2MM payload FIFO를 포함한다.
MM2S FIFO는 68 bit, S2MM FIFO는 72 bit라서 각각 RAMB36 하나와 남은 2 bit용
RAMB18 하나를 사용한다. 따라서 accelerator의 6 RAMB36에
\[
2\times\mathrm{RAMB36}+2\times0.5\,\mathrm{RAMB18}=3\ \text{tiles}
\]
가 더해져 9 tile이 된다. 이것은 AXI4-Stream protocol 자체나 NTT arithmetic의
비용이 아니라 해당 DMA FIFO configuration의 data-movement 비용이다.

AXI BRAM Controller configuration은 이 payload FIFO 없이 coefficient memory에
접근하며, 현재 post-route 결과에서 controller 자체는 BRAM을 사용하지 않는다.
complete system의 area 차이를 core arithmetic 개선으로 해석하면 안 된다.

\begin{checklistbox}
\begin{itemize}
  \item layer당 issue, fill, drain, transition cycle을 분리해 센다.
  \item logical schedule와 physical mapping을 독립 module/test로 둔다.
  \item ready/valid throughput과 total latency를 별도로 측정한다.
  \item in-flight response까지 FIFO credit에 포함한다.
  \item interface 비교는 같은 device, clock, tool, implementation directive로 한다.
\end{itemize}
\end{checklistbox}

